This chapter provides an introduction to the topic of the present project work.

This work is being written by Timon Burkard and Matthias Schär. The supervising lecturers are Prof. Guido Keel and
Michael Lehman.

\subsection{Topic}

The work deals with the development of demonstrator electronics for the visualization of optical \acrfull{tof} for
distance measurement.

A measurement chip, the TDC7200 from Texas Instruments, is put into operation. This is done in two steps: first, purely
electrical measurements are to be made possible, for example to determine the length of connected cables. Furthermore,
the optoelectronics needed for optical distance measurement are to be located on the demonstrator.

\subsection{Objective}

The students have set themselves the goal of developing electronics that can be used to carry out the experiments. The
control of the measuring equipment and the design of the optoelectronics should also be largely carried out independently.

The aim is to create a demonstrator that can be used with the TDC7200 to determine different cable lengths of determine
the length of copper cables. After this goal has been achieved, it should also be possible to carry out optical
measurements with the same demonstrator. The students have set themselves the goal of ensuring that the measurements
with a wall as a target work up to a distance of 10 meters. Furthermore, a spatial resolution of about 30 centimeters
should be achieved.

\subsection{Structure}

The following is a brief overview of the chapters in the rest of this paper:

\begin{itemize}
    \item \textbf{Chapter~\ref{sec:theory} - \nameref{sec:theory}:} Deals with the theoretical foundations that were
          developed for the demonstrator.

    \item \textbf{Chapter~\ref{sec:concept} - \nameref{sec:concept}:} Explains the rough concepts and challenges that
          need to be overcome. In particular, the electronics to be developed and the time requirements for the
          \acrshort{tdc} are recorded.

    \item \textbf{Chapter~\ref{sec:simulations} - \nameref{sec:simulations}:} Contains information about the
          simulations that are carried out before the design but also during the commissioning of the demonstrator.

    \item \textbf{Chapter~\ref{sec:realization} - \nameref{sec:realization}:} Explains the realization. Circuits and
          the manufacturing process are presented here. It also explains what software development is necessary to
          realize such a demonstrator.

    \item \textbf{Chapter~\ref{sec:results} - \nameref{sec:results}:} Documents all measurements carried out and their
          results. It shows to what extent the initial objective has been achieved.

    \item \textbf{Chapter~\ref{sec:conclusion} - \nameref{sec:conclusion}:} Captures the students' final thoughts and
          summarizes the work. It also provides an outlook on possible extensions and improvements.
\end{itemize}
